% CLAUDE REPORT::SWITCH-0001::2026-11-14T09:00-08:00 (PST)
\documentclass[11pt]{article}
\usepackage{amsmath,amssymb,amsthm}
\usepackage[margin=1in]{geometry}
\usepackage{url}
\newtheorem{theorem}{Theorem}\newtheorem{lemma}{Lemma}\newtheorem{corollary}{Corollary}
\title{CLAUDE-SWITCH-0001: Adaptive strategy switching in the worst-case\\ minimum-derivation objective $\overline W_R(N)$}
\author{Claude (scheduled worst-case-derivation routine)}
\date{14 November 2026, 09:00 PST}
\begin{document}\maketitle

\noindent\textbf{Status summary.} Lemma~1 PROVED/EXACT (known technique: decision-tree restriction argument). Theorem~2 PROVED/EXACT under its stated hypotheses. The instantiation in the historical mixed calculus is CONDITIONAL on a guarded-row rule. Finite checks are COMPUTATIONAL. The following remain \textbf{OPEN}: the boxed bridge $\mathrm{SAT}\in\mathrm P\Rightarrow\overline W_R(N)\le\mathrm{poly}(N)$, any superpolynomial $\overline W_R$ for a switching-closed $R$, and P versus NP.

\section{Target}
The user's shared objective is
\[
\overline W_R(N)=\max_{F\in\mathrm{UNSAT},\ |\langle F\rangle|\le N}\ \min_{\pi:F\vdash_R\bot}\mathrm{cost}_R(\pi),
\]
where the best derivation may be chosen for free. This run addresses priority~4: do proposed hard instances survive \emph{adaptive strategy switching}? Here that means mixing inference systems inside one derivation, rather than picking the best single system for each instance. Minimizing over derivations is a valid lower-bound device \emph{for the calculus over which one minimizes}. The question is which calculus must be fixed.

\section{Guarded composition}
For UNSAT CNFs $F_1,F_2$ and a fresh variable $z$, let
\[
G(F_1,F_2)=\{C\vee z: C\in F_1\}\cup\{C\vee\neg z: C\in F_2\}.
\]
$G$ is UNSAT, with $G|_{z=0}=F_1$ and $G|_{z=1}=F_2$. In CNF form, it says ``first decide which strategy applies''.

\begin{lemma}[exact additivity, tree-like Resolution]
Let $L(F)$ be the minimum number of leaves in a tree-like Resolution refutation of $F$ (with weakening). Equivalently, $L(F)$ is the minimum number of leaves in a decision tree whose leaves each falsify a clause of $F$. Then $L(G(F_1,F_2))=L(F_1)+L(F_2)$. Disjointness of variables is not needed.
\end{lemma}
\begin{proof}
($\le$) Query $z$ first, then attach optimal trees for $F_1$ and $F_2$.

($\ge$) Let $D$ be any decision tree for $G$. Every clause of $G$ contains a $z$-literal, so a leaf whose path does not assign $z$ falsifies no clause of $G$. Hence every leaf assigns $z$. Write $\Lambda_b$ for the leaves whose path sets $z=b$; the sets $\Lambda_0,\Lambda_1$ partition the leaves. Restricting $D$ by $z=0$ deletes the $z=1$ subtrees and leaves a decision tree with leaf set exactly $\Lambda_0$. Each leaf in $\Lambda_0$ falsifies some $C\vee z$, hence falsifies $C\in F_1$, so this is a decision tree for $F_1$ and $|\Lambda_0|\ge L(F_1)$. Symmetrically, $|\Lambda_1|\ge L(F_2)$.
\end{proof}
Within one restriction-closed system, the step ``decide which strategy'' therefore costs exactly the sum: it gives neither amplification nor savings. So hard instances for $\overline W_R$ must be irreducible under this composition, which agrees with the corpus's MUS worst-case quotient (LSPV-1/350).

\begin{theorem}[switching gap]
Let $R_1,R_2$ be clause-line refutation systems whose cost does not increase under restrictions. Let $R_1\sqcup R_2$ be the system that permits each step from either one. Assume (H-lift): every $R_i$ refutation of $F$ of cost $c$ lifts to an $R_1\sqcup R_2$ derivation of the unit $z$ (respectively $\neg z$) from the guarded copy, at cost $\le c+O(1)$. Then
\[
W_{R_i}(G)\ \ge\ \max\{W_{R_i}(F_1),W_{R_i}(F_2)\}\quad(i=1,2),\qquad
W_{R_1\sqcup R_2}(G)\ \le\ W_{R_1}(F_1)+W_{R_2}(F_2)+O(1).
\]
\end{theorem}
\begin{proof}
The lower bound follows from restriction closure, since $G|_{z=0}=F_1$ and $G|_{z=1}=F_2$. For the upper bound, lift the two refutations to derivations of $z$ and $\neg z$, then resolve once.
\end{proof}
\begin{corollary}
Suppose $F_1$ is easy for $R_1$ and hard for $R_2$, and $F_2$ is easy for $R_2$ and hard for $R_1$. Then on $G$ the best \emph{single} system is expensive, $\min_i W_{R_i}(G)\ge\min_i\max_j W_{R_i}(F_j)$, while the switching calculus is cheap. Consequently, checking separately that a family survives Resolution and that it survives algebraic reasoning does \textbf{not} certify a lower bound against an algorithm that switches strategies. The objective must fix one switching-closed $R$. If $R_1$ p-simulates $R_2$, the gap disappears.
\end{corollary}

\paragraph{Instantiation (CONDITIONAL).} Take $R_1=$ Resolution and $R_2=$ GF(2) row-XOR. Row-XOR is incomplete, so $W_{R_2}=\infty$ on non-affine CNFs. Let $F_1$ be any non-affine CNF with a short Resolution refutation, and let $F_2$ be an odd Tseitin formula on a constant-degree expander. Resolution needs exponential size on $F_2$, $W_{\mathrm{Res}}(F_2)=2^{\Omega(n)}$ (Urquhart 1987), while summing all vertex equations gives a linear GF(2) refutation. Hence $\min(W_{\mathrm{Res}}(G),W_{\mathrm{XOR}}(G))=2^{\Omega(n)}$, but $W_{\mathrm{Res}\sqcup\mathrm{XOR}}(G)=O(|G|)$.

For XOR, H-lift needs a sound guarded-row rule, $(z\vee A),(z\vee B)\vdash z\vee(A\oplus B)$. Whether the historical $R_{\mathrm{mix341}}$ grammar natively supports it is unauthenticated, as LSPV-1/592 already flags for other native rules. The instantiation is consistent with LSPV-1/540's observation that Urquhart's family fails against the mixed calculus ($4h-1$ nodes).

\section{Literature clue for this obstacle}
Up to simulation, the natural switching closure of priority~4's two tests (Resolution plus Gaussian elimination) is Res($\oplus$)/Res(lin). Exponential lower bounds are known for \emph{tree-like} Res($\oplus$), for example for the pigeonhole principle (Itsykson--Sokolov; extended to other fields by Part--Tzameret), and recent work targets \emph{regular} ResLin (ECCC TR24-022). According to the sources consulted this run, however, \textbf{no superpolynomial lower bound is known for dag-like Res(lin)}. So the smallest switching-closed $R$ that passes both of the user's tests already sits at a recognized open frontier of proof complexity.

Adding circuit compression (extension variables) gives Extended Resolution, which is p-equivalent to Extended Frege (EF). Cook (1975) showed that EF is the weakest system that is polynomially bounded on the propositional translations of PV-provable coNP statements. No superpolynomial lower bounds for EF are known.

\begin{corollary}[bridge ladder, KNOWN THEOREM assembly]
\begin{enumerate}
\item[(a)] Let $R$ be a system with an already proved superpolynomial $\overline W_R$, such as Resolution (Haken 1985; Urquhart 1987). For such $R$, the boxed bridge is equivalent to $\mathrm{SAT}\notin\mathrm P$, so it is no weaker than the target (compare REFEREE296).
\item[(b)] For $R=$ EF, the bridge holds when the hypothetical polynomial algorithm's UNSAT-soundness is PV-provable (Cook 1975). The exact gap is removing ``PV-provable''. LSPV-1/357 and LSPV-1/415 show this gap cannot be closed by uniform reflection.
\item[(c)] For dag-like Res(lin) and every stronger switching-closed $R$, both the bridge and the lower bound are open.
\end{enumerate}
\end{corollary}

\section{Executed tests (COMPUTATIONAL)}
The program \path{experiments/claude/claude_guarded_switch001.py} (git blob \texttt{9161e54084539bc662bb4190d1304b930d45d8d5}) ran natively in 5.07\,s as a single chunk. Raw output is in \path{results/raw/claude/2026-11-14_0900_CLAUDE_SWITCH0001_RAW.json} (blob \texttt{09b60bb5ae8f0a8404bd1d419625c003e3855b3d}).
\begin{itemize}
\item \textbf{Random pairs.} 80 random UNSAT pairs with at most $5+4+1$ variables. For each pair, the program checked that $G$ is UNSAT, that the $z=0$ restriction identity holds, and a SAT control ($G$ with the $F_2$ part removed is satisfiable). The exact equality $L(G)=L(F_1)+L(F_2)$ held for all 80 pairs, with 0 failures.
\item \textbf{Named instances.} $L(\mathrm{PHP}^3_2)=11$ and $L(\text{odd Tseitin on }K_3)=6$; their composite has $L(G)=17$, as the equality predicts. Odd Tseitin on $K_4$ needs $L=28$ leaves, compared with 3 row-XORs over GF(2).
\end{itemize}
Finite data do not bear on asymptotic lower bounds, and leaf count is not machine time or space.

\section{CORPUS REVIEWED}
\begin{itemize}
\item AGENTS.md, all directives. This includes the 2026-11-09 knot handoff, which is not acted on here.
\item The TMR replacement directive \path{research/coordination/2026-11-13_1249_TM_RESOURCE_AGENT_REPLACEMENT.md} (blob 47459817).
\item \path{research/coordination/2026-11-14_0819_USER_CLAUDE_SEMANTIC_COMPOSITION_REVIEW.md}.
\item The TMR-1 and TMR-5 handoffs, and the directory listings of TMR-1 through TMR-5.
\item The $\overline W_R$ and B1--B3 passages of root \path{main.tex}: REFEREE295 and 296, and LSPV-1/350, /351, /353, /357, /540 and /592.
\item The opening sections of LSPV-1/406 and LSPV-1/415.
\end{itemize}
Not fully reread this run: proof1.txt, proof2.txt and issue \#3. No tracker change is made.

\section{DUPLICATION / CONTRADICTION AUDIT}
The reflection and EF bridges (LSPV-1/357, /406, /415, /592 and REFEREE296) are credited, not reproved. The Urquhart-versus-mixed failure duplicates LSPV-1/540 and is re-derived only as an instance of Theorem~2. Searches of the corpus found neither the guarded additivity equality nor a stated switching-gap requirement for $\overline W_R$; these are this run's new increment. No contradiction was found.

\section{WHAT CHANGED RELATIVE TO ROOT MAIN}
Root \path{main.tex} is \textbf{not edited}. TMR-5 is its current synthesis owner, and this routine has no owned region. Proposed for TMR-5 review: add the rule ``hard families must be certified against one switching-closed $R$'' to the falsification ledger.

\section{ADOPTED / REJECTED / RETIRED PRIOR COMPONENTS}
Adopted: the max--min objective with free choice of derivation, which is valid per calculus. Rejected as a lower-bound certificate: ``the family survives each system separately''. Nothing retired.

\section{STALE OR OVERSTATED CLAIMS FOUND}
None identified in the passages read.

\section{Post-quantum ledger}
Item 6 (collision extraction): guarded composition resembles a binding-style test. This is an ANALOGY only. Items 1--3: FAILED TRANSFER in this run.

\section{ONE NEXT TARGET}
Fix the minimal switching-closed calculus $R_\star=\mathrm{Res}(\oplus)$ with charged guarded rows. Then test whether any corpus candidate family, including the TMR sign-stream $U_r$, has polynomial-size $R_\star$ refutations. Only a candidate with no known short $R_\star$ refutation is worth carrying into $\overline W_R$.

\paragraph{References.} S. Cook and R. Reckhow, J. Symb. Logic 44 (1979). S. Cook, STOC 1975 (PV). A. Haken, TCS 39 (1985). A. Urquhart, JACM 34 (1987). D. Itsykson and D. Sokolov, MFCS 2014 / APAL 2020. F. Part and I. Tzameret, ITCS 2020 / ECCC TR20-034. ECCC TR24-022. J. Pich and R. Santhanam, ECCC TR23-199.

\textbf{P versus NP remains OPEN.}
\end{document}
% END_OF_VERIFIED_REPORT::CLAUDE::SWITCH-0001
